\section{Related Work}
\label{sec:related}

Our work bridges LLM code repair and feedback mechanism design. We review each area, highlighting the gap our matched-content design addresses.

\subsection{LLM Self-Repair}

Iterative self-repair has become the dominant paradigm for improving LLM code generation beyond single-shot accuracy. \citet{olausson2023selfrepair} demonstrated that feeding execution feedback (test results, stack traces) back to the model enables 12--17\% relative improvement on HumanEval. \citet{arimbur2026tries} showed that most gains concentrate in the first 2--3 repair iterations, with diminishing returns thereafter.

Several frameworks operationalize self-repair. CodeRL \citep{le2022coderl} integrates reinforcement learning with execution feedback. ThinkRepair (ISSTA 2024) uses self-directed debugging on Defects4J. More recent work explores multi-agent architectures: CodeCoR (2025) achieves 77.13\% pass@1 on HumanEval/MBPP through agent collaboration.

These approaches share a common assumption: feedback content determines repair quality. The question of feedback \emph{ordering}---given the same content---remains unexplored.

\subsection{Static Analysis Feedback}

Static analysis provides complementary signals to execution feedback. \citet{jain2024llmcleaning} showed that LLM-assisted code cleaning improves code quality without runtime signals. \citet{blyth2025static} demonstrated that Pylint feedback reduces security issues from >40\% to 13\% and reliability warnings from >50\% to 11\%.

AutoSafeCoder \citep{nunez2024autosafecoder} combines static analysis with fuzz testing in a multi-agent framework, achieving 13\% vulnerability reduction. \citet{dolcetti2024helping} integrated testing and static analysis feedback for safety improvements.

However, studies comparing static and execution feedback face a fundamental confound: they necessarily vary information volume alongside feedback type. Cascaded approaches (static + execution) provide more information than single-type approaches. Our prior work (h-c1) achieved 16.18\% improvement with cascaded feedback, but whether this stems from \emph{more} information or \emph{ordered} information was unclear.

\subsection{Feedback Presentation Effects}

Research on prompt engineering and positional effects suggests LLMs are sensitive to presentation order. Recency effects in transformer attention mean later tokens receive disproportionate weight. This raises an alternative hypothesis: static-first might help not because it structures the repair process, but simply because execution feedback appears last (terminal position bias).

FeedbackEval \citep{dai2025feedbackeval} benchmarked feedback-driven code repair across feedback types, finding mixed feedback yields 63.6\% repair success. However, ordering effects were not systematically tested.

\subsection{Our Contribution}

Unlike prior work, we control for information volume by ensuring byte-identical feedback across conditions. This matched-content design isolates the ordering effect from the content effect. We test both the existence of an ordering effect and its mechanism---distinguishing scaffolding (coarse-to-fine repair) from recency bias (terminal position advantage).
